\section{Actual first-step NC histories in coordinate models}
\label{sec:hironaka-normal-crossings-first-step-chart-history}

Let $H$ be an actual NC carrier in a complex $n$-manifold, with positive
highest branch count $r$, and let $C=\{x:r\leq\nu_H(x)\}$. We retain
the unchanged original datum $M_H=(X,X,\varnothing,I_H,1)$ and its genuine
first blowup along $C$. Coordinate models below are actual blowups of
the coordinate centre, not just present-day descriptions of strict transforms.
Model comparison points remain actual images of the original blown-up
points. No fixed orbit representative or finite whole-value image is asserted.

\begin{lemma}[A genuine one-step history isomorphism from actual blowups]
\label{hir:general-nc-one-step-history-universal-comparison}
\leanok
\lean{Hironaka.HironakaNormalCrossingsFirstStepChartHistory.exists_oneStep_isTowerIso}
An actual biholomorphism between two bases that identifies their whole
centres lifts to a genuine isomorphism between any actual blowups of
those centres. The lifts commute exactly with the projections, and the
result is an isomorphism of the entire one-step towers.
\end{lemma}
\begin{proof}
Compose the second blowup with the inverse base biholomorphism. Its
actual blowup universal property is the blowup property over the first
base along the first centre. The two actual blowups over that same base
therefore factor uniquely through one another. Their composites are
holomorphic self-maps over the projection; uniqueness in the universal
property makes both composites the identity. These comparison maps give
an actual biholomorphism of the blown-up manifolds. Apply the base
biholomorphism to the exact projection equation to obtain exact
commutation over the original base map. Together with the assumed whole
centre identity, the reflexive base history extends to the asserted
one-step history isomorphism. No lifted map or history-isomorphism
certificate is supplied.
\end{proof}

\begin{lemma}[Literal whole carrier and centre identities in an actual NC chart]
\label{hir:general-nc-first-step-whole-chart-data}
\leanok
\lean{Hironaka.HironakaNormalCrossingsFirstStepChartHistory.chart_preimage_centre,Hironaka.HironakaNormalCrossingsFirstStepChartHistory.chart_preimage_carrier}
In an actual highest-stratum chart, the biholomorphism from its true
open domain to its lifted open coordinate target pulls the coordinate
centre back to the entire trace of $C$ and the selected coordinate union
back to the entire trace of $H$.
\uses{hir:general-nc-highest-stratum-branch-adapted-coordinates}
\end{lemma}
\begin{proof}
At every point of the actual chart domain, the chart's centre equation
identifies the intrinsic highest-count condition with vanishing of all
last $r$ coordinates. This is precisely membership in the actual coordinate
centre. Its carrier equation identifies membership in the whole original
carrier with vanishing of at least one of those coordinates. These
pointwise equivalences prove the two literal preimage set identities.
Both statements hold on the whole open domain, including points outside
the original centre or carrier.
\end{proof}

\begin{theorem}[An internally produced first-step coordinate history]
\label{hir:general-nc-actual-first-step-coordinate-history}
\leanok
\lean{Hironaka.HironakaNormalCrossingsFirstStepChartHistory.exists_highest_chart_history,Hironaka.HironakaNormalCrossingsFirstStepChartHistory.HighestChartHistory.pullback_original}
Given an actual blowup along $C$, a point $x\in C$, and any fixed genuine
coordinate-centre blowup, the original NC carrier internally produces an
actual centred highest-stratum chart and an isomorphism from the original
one-step history restricted to its domain to the coordinate blowup history
restricted to the chart target. The original marked data are identified
in full with the original coordinate-carrier datum on that target.
\uses{hir:general-nc-one-step-history-universal-comparison,hir:general-nc-first-step-whole-chart-data,hir:general-nc-generic-initial-value-flat-chart}
\end{theorem}
\begin{proof}
The highest-stratum chart producer gives an actual centred flattening
at $x$, with the whole centre and carrier equations of the preceding
lemma. Restrict each of the two actual blowups to the corresponding
actual open base. Genuine blowup restriction gives the required two
universal properties, and the preceding whole centre identity identifies
their centres under the actual chart biholomorphism. Apply the universal
comparison lemma to construct an actual biholomorphism of their entire
inverse-image opens and a genuine one-step tower isomorphism.

For the marked records, the whole carrier identity and the actual chart
and its inverse identify holomorphic germs vanishing on the two whole
traces. Thus their complete vanishing ideals agree after actual pullback.
The remaining fields agree literally: full ambient carrier, empty old
boundary and marking one. This identifies the whole original records,
without replacing them by present strict-transform data. All maps and
inverse-image domains are those of the actual first blowups.
\end{proof}

\begin{theorem}[Whole ambient-dimensional canonical values on the actual first histories]
\label{hir:general-nc-first-step-whole-canonical-chart-value}
\leanok
\lean{Hironaka.HironakaNormalCrossingsFirstStepChartHistory.exists_highest_chart_history_canonical_value}
Assume $n>0$. In the preceding construction, the entire first-step value
of $\mathrm{canonicalInvariant}(n)$ at each actual transformed cosupport
point equals its value on the literal original coordinate-carrier datum
along the genuine restricted coordinate blowup history at the image
point. All chronological and born-boundary information is retained by
the actual one-step history isomorphism.
\uses{hir:general-nc-actual-first-step-coordinate-history,hir:general-nc-ambient-canonical-initial-full-maximum}
\end{theorem}
\begin{proof}
The actual original full-maximum theorem identifies $C$ with the entire
initial maximum for the exact ambient-dimensional canonical invariant.
The given point of $C$ supplies attainment, so the actual first history
is canonical without a supplied history certificate. Restricting to the
chosen chart domain preserves the initial whole maximum: the domain
meets $C$, and the proved reflexive restriction law preserves every whole
value and its cosupport. Hence the restricted first step is canonical
for the same original pulled-back datum.

Apply the invariant's genuine whole tower-isomorphism law to the inverse
of the constructed one-step history isomorphism. Its right-hand original
restricted history is canonical as just proved. The whole original-record
identity turns the left-hand datum into the literal coordinate-carrier
datum on the actual open target. The original cosupport restriction
identity supplies precisely the required point hypothesis. Finally apply
the invariant's proved open restriction law to the global original first
history. The composite of these two whole-value equalities gives the
claim. The model point is the actual image under the constructed lift;
no assertion that it is a fixed representative, or that finitely many
whole values occur, follows without an additional orbit construction.
\end{proof}
